\section{Discussion}
\label{sec:discussion}

We studied a top-down preparation workflow that separates ideal circuit
construction, gate removal, and noisy refinement. Exact single-gate
deletion scores identify the effect of each removal on the current ideal
target fidelity. Frozen Pauli histories make the sampled refinement
repeatable, and empirical weighting preserves the contributions of
duplicate no-jump paths. Independent validation selects the circuit
before exact evaluation of the same channel. These operations give a
specified and auditable preparation procedure.

The distinctions between its stages matter. Exact deletion damage for
one gate does not bound the combined damage from simultaneous removals.
Restoring the best noiseless parameters protects the returned ideal
checkpoint, while the noisy phase retains its final iterate. Its
cross-trajectory signal is not identified with the sampled-loss gradient.
The noiseless adjoint interpretation and quasi-Newton safeguards therefore
do not provide a convergence or monotonicity guarantee for the complete
noisy workflow. The numerical results evaluate the implemented procedure
with these limitations intact.

Under logical-gate noise, the method combines the highest selected fidelity
on five of nine targets with Pareto efficiency in fidelity and classical
cost on seven. The six real targets give the strongest trade-offs against
RobustState. Fontana remains an inexpensive alternative on that subset.
Complex Haar targets require more top-down construction work, and
RobustState dominates two of the three instances. The internal controls
support the usefulness of the full workflow at the tested settings, but
their differing policies prevent attributing the observed gains to a
single component.

Noise placement is a substantive part of this conclusion. Logical
two-qubit primitives do not have equal physical decomposition costs.
The native pilot gives a weaker top-down comparison on Gaussian and
Haar 1, while retaining a classical-cost trade-off on critical Ising.
The same-schedule parameter diagnostic also gives little benefit from
switching the refinement to native noise at the selected endpoints.
Thus the favorable logical results should be interpreted together with
the native results rather than extrapolated to a common physical gate
implementation.

RobustState's refinement benefit also depends on the noise mechanism.
Its authors attribute the hardware gains primarily to coherent-error
correction and report only marginal fidelity improvement in their
incoherent-noise simulator~\cite[Sec.~IV.C]{RobustState2024}.
Our Pauli-noise benchmark excludes systematic coherent gate errors and
therefore probes a related regime in which additional benefit from
parameter refinement may be limited. This qualifies the expected
refinement gain, rather than the achievable preparation fidelity:
RobustState remains competitive on our complex Haar targets.
The present comparison does not assess its reported advantage in
correcting coherent hardware errors.

Separate twelve-qubit diagnostics indicate a substantial cost for noisy
IQSP in our benchmark implementation. An early depth-four Gaussian
checkpoint required $505\,\mathrm{s}$ for one full fidelity-gradient
evaluation with our default compression on one Apple M1 Pro thread; peak
physical footprint was $4.76\,\mathrm{GiB}$.
% Provenance and compression settings: data/q12_iqsp_resource_diagnostic.json.
The published H2 demonstration uses noiseless one-qubit gates and
angle-dependent, factorized single-qubit depolarization after native $ZZ$
rotations~\cite{Ballarin2025}. Our benchmark additionally includes joint
two-qubit Pauli errors, requiring joint noise contractions and compression
in the locally purified implementation. Our trajectory refinement instead
propagates sampled pure states without purification indices, at the cost
of sampling error. This suggests a potential computational advantage for
parameter-independent Pauli noise. These checkpoint diagnostics neither
isolate noise structure as the cause of the observed cost nor establish
superior end-to-end performance or scaling. IQSP itself is not restricted
to the noise model or circuits used in its demonstration.

The main comparison is limited to a development-used fixture, six qubits, one
optimizer seed, one synthetic Pauli profile in the all-target comparison,
and finite, unequal search budgets. The exact-density feedback supplied
to several comparators is a simulator adaptation. It does not measure
their hardware sampling costs or their response to coherent device
errors. The dense implementation also does not establish the scaling of
a tensor-network realization. In particular, the retained-history cross
update contains contractions over pairs of trajectories; duplicate
compression alone does not imply linear scaling with the raw sample
count.

Further evaluation should vary optimizer seeds and use new target
instances, extend the matched native study, and test budget sensitivity.
Matched component ablations would be needed to separate the roles of
pruning, construction refinement, and the noisy update. Mathematical
analysis of the cross-trajectory signal is a separate question from
implementation correctness. Within the measured setting, top-down
pruning and refinement provide a useful preparation option whose value
depends on the target, circuit representation, and available classical
computation.
